% !TeX root = ../main.tex
% ============================================
% 第三章 系统需求分析与总体设计
% ============================================

\chapter{系统需求分析与总体设计}

本章在前述理论基础之上，进一步说明平台面向的核心需求、总体架构、数据模型与关键模块设计。与第4章侧重算法方法、第5章侧重工程实现不同，本章主要回答“系统为何这样划分模块、各模块之间如何协同”这一设计问题，从而为后续的实现与实验分析提供结构化支撑。

\section{系统需求分析}

\subsection{功能需求分析}

结合平台的目标用户与实际业务链路，系统功能需求可以归纳为三个方面。

第一，系统需要支持多模态健康数据的采集、整理与管理。用户既可以直接录入临床指标与基本画像，也可以上传体检报告，由 OCR 模块辅助提取关键字段；对于可持续采集的行为数据与设备数据，系统还应支持按时间维度持续更新用户健康画像。该需求对应前端健康档案、报告上传与时间线展示等功能，也是后续风险分析的输入基础。

第二，系统需要提供面向个体的风险评估与干预推演能力。平台应能够在统一分析入口下整合临床、遗传与行为信息，输出结构化风险结果，并进一步支持未来风险模拟与干预情景比较，帮助用户理解不同生活方式调整对风险变化的影响。

第三，系统需要提供可解释的智能健康问答能力。问答模块不仅要支持连续对话，还要能够在回答中结合用户画像、历史分析结果与检索到的医学资料，输出带有证据标记和必要提示信息的响应；对于家庭协同场景，系统还应支持家庭账户关联，便于在授权前提下完成代录入与健康关注。

除面向普通用户的核心功能外，系统还需要提供基础管理能力，用于数据导入、模型训练任务触发、知识库维护和系统状态查看。这些功能主要服务于平台运行维护，不作为正文中的主叙事重点，但在架构设计中仍需预留对应接口与角色边界。

\subsection{非功能需求分析}

在非功能层面，平台主要面临安全隐私、交互响应与可维护性三方面约束。

首先，平台处理的健康档案、体检文档和遗传数据均具有较高敏感性，因此系统需要具备用户鉴权、权限隔离与敏感字段保护能力。结合现有实现，用户访问依赖登录令牌进行校验，遗传信息采用加密存储方式保存，问答链路中的治理摘要与审计留痕则以结构化形式写入数据库，以便后续追踪与回放。

其次，平台需要兼顾分析任务与日常交互的响应体验。风险评估、文档解析与问答服务均通过统一接口向前端提供结果，而较长时间的数据导入或训练任务则采用后台异步方式执行，以避免阻塞在线交互流程。

最后，平台需要具备一定的扩展性与可维护性。考虑到健康指标字段、知识资料和问答治理逻辑会随着平台迭代而逐步扩展，系统设计上采用前后端分离、服务分层与半结构化数据存储的组合方式，以降低后续功能增量对整体架构的扰动。

\section{系统总体架构设计}

\subsection{平台总体架构}

结合仓库中的前端、后端与部署配置，平台总体上采用“前端交互层、后端服务层、数据与模型支撑层”三层结构。

前端交互层基于 Vue 构建，负责健康档案录入、风险结果展示、时间线浏览以及 Dr.~AI 问答等界面组织，并通过统一路由区分普通用户与管理员入口。后端服务层基于 FastAPI 构建，承担用户认证、综合分析、文档解析、会话管理、问答治理与后台任务调度等职责，是整个平台的核心业务承载层。数据与模型支撑层则由关系型数据存储、Redis 缓存以及风险评估、OCR、知识检索等运行时组件共同组成，用于支撑分析结果生成与状态持久化。

在这一结构中，Agent 健康问答模块不作为独立系统存在，而是位于后端服务层内部，与用户画像、风险分析结果、知识检索和审计记录协同工作。这样既能复用平台已有的健康数据上下文，也便于将问答响应与风险评估结果放在统一交互入口下展示。

\begin{figure}[ht]
    \centering
    % [图 3-1 平台总体架构图]
    \caption{多模态健康管理平台总体架构示意}
    \label{fig:overall_architecture}
\end{figure}

\subsection{核心业务流程设计}

从核心业务流程看，平台的运行链路可以概括为“数据输入、统一分析、结果反馈、持续问答”四个环节。

首先，用户通过健康档案页面填写临床画像，或上传体检报告触发 OCR 解析；必要时，系统还可以补充遗传数据与行为数据，形成较为完整的分析输入。随后，后端在统一分析入口下对多源数据进行整理与融合，生成结构化风险结果、分析上下文及可供后续展示的解释信息。分析完成后，前端以仪表盘、风险卡片、时间线等方式呈现结果，并在需要时继续调用未来风险模拟或干预模拟接口，向用户展示不同场景下的变化趋势。

在反馈环节中，Dr.~AI 问答模块进一步利用当前用户画像、历史会话与知识检索结果组织回答。当用户提出健康咨询、报告解读或趋势回顾类问题时，系统根据问题类型选择受限工具与证据来源，输出包含回答文本、证据标签和必要提示的结构化结果，从而形成“分析结果驱动问答、问答反过来解释分析结果”的协同流程。

\begin{figure}[ht]
    \centering
    % [图 3-2 核心业务流程图]
    \caption{平台核心业务流程示意}
    \label{fig:core_flow}
\end{figure}

\section{核心数据模型设计}

\subsection{关键实体与关系}

平台的数据模型围绕用户主体、健康观测、问答会话与治理留痕四类对象展开。用户账户由 \texttt{User} 实体负责认证与角色标识，\texttt{UserProfile} 与之构成一对一关系，用于保存临床画像、遗传信息摘要和扩展健康字段。围绕用户主体，\texttt{HealthRecord} 用于保存历史观测快照，\texttt{MedicalDocument} 用于保存上传文档及其 OCR 处理结果，二者共同构成健康数据的历史积累基础。

在交互层面，\texttt{ChatConversation} 与 \texttt{ChatMessage} 用于组织对话会话及逐条消息内容，从而支持连续问答与历史回溯。对于家庭协同场景，\texttt{FamilyLink} 记录授权关系，使系统能够在明确权限边界的前提下支持家庭成员之间的健康关注与代录入。

围绕 Agent 问答治理，系统进一步引入 \texttt{AgentAuditEvent} 与 \texttt{AgentAnswerReplay} 两类实体，用于保存车道判定、回答模式、证据充分性、工具使用摘要以及回答回放所需的关键快照信息。这部分数据不直接面向用户展示，但为后续审计、回归验证和问题排查提供了结构化依据。

\begin{figure}[ht]
    \centering
    % [图 3-3 核心实体关系图]
    \caption{平台核心数据实体关系示意}
    \label{fig:core_er_model}
\end{figure}

\subsection{多模态数据组织方式}

为兼顾结构稳定性与后续扩展需求，平台采用“核心字段结构化存储，扩展信息半结构化保存”的数据组织方式。对于年龄、血糖、糖化血红蛋白、血压等稳定临床指标，系统将其作为 \texttt{UserProfile} 的显式字段保存，便于后续分析服务直接读取。对于随时间变化的观测值及风险快照，系统通过 \texttt{HealthRecord} 进行按时间记录，从而支撑历史趋势展示与阶段性比较。

文档类输入则通过 \texttt{MedicalDocument} 承载文件路径、上传时间与 OCR 处理状态，使原始文件保存、解析状态追踪与结构化导入可以在同一对象下管理。对话类数据由会话实体和消息实体分别保存，以便同时支持最近消息窗口构建、历史检索和前端会话列表展示。对于问答治理中产生的证据面板、上下文预算摘要和工具调用摘要，系统使用 JSON 字段保存结构化扩展信息，以平衡模式稳定性与表达灵活性。

需要指出的是，平台对遗传相关信息采取了独立加密字段存储方式，而不是将其完全混入普通文本字段中；这既符合健康数据的敏感性要求，也与当前代码中的属性访问和加解密处理逻辑保持一致。

\section{关键模块设计}

\subsection{风险评估模块设计}

风险评估模块是平台分析链路的核心入口。其主要任务是接收用户临床画像、遗传数据和部分行为信息，在统一分析接口下完成特征整理、风险估计与结果组织，并输出可供前端展示的风险报告和分析上下文。该模块在设计上强调“统一入口、统一输出”，从而避免不同数据源分别分析后难以汇总的问题。

从系统协同关系看，风险评估模块一方面与用户画像、历史记录和文档解析结果相连，负责消费结构化健康数据；另一方面又向仪表盘展示、时间线展示和问答模块提供统一的风险结果支撑。因此，它在整个平台中处于“由输入走向决策”的桥接位置，是第4章算法设计在工程层面的主要承载对象之一。

\subsection{干预推演模块设计}

干预推演模块建立在风险评估结果之上，主要提供未来风险模拟和干预情景比较两类能力。系统通过独立接口接收用户当前画像及目标干预参数，并输出自然演化场景或干预后场景下的风险变化结果。相较于直接给出静态建议，这种设计更强调“可比较”和“可解释”，有助于用户理解不同调整方案可能带来的风险变化方向。

在平台链路中，干预推演模块与风险评估模块共享核心输入字段，但其输出更侧重面向用户的解释与比较，因此会直接服务于前端中的趋势展示、时间线和干预建议交互。第4章将进一步说明其方法基础，第5章则对应其接口与服务实现。

\subsection{Agent健康问答模块设计}

Agent 健康问答模块承担平台中的智能交互功能。该模块接收用户自然语言问题，在会话上下文的基础上结合用户画像、历史分析结果、文档摘要和知识检索结果生成回答，并根据问题类型调用受限工具集完成补充查询。与普通问答不同，该模块的设计重点不在开放式生成，而在于受控调用、证据约束和结果留痕。

从数据流角度看，问答模块以上一轮分析结果和当前会话状态为输入，以结构化回答结果为输出。输出除了回答文本外，还包括来源引用、证据标签、建议卡片以及必要时的人工接管提示。与此同时，系统将本轮问答过程中的关键治理摘要写入审计与回放对象，以便后续进行行为验证和责任追踪。由于这部分内容与实际运行链和服务编排紧密相关，其具体实现将在第5章中集中展开，而本章仅保留其模块定位与职责说明。

\section{本章小结}

本章从需求、架构、数据模型与关键模块四个层面说明了平台的总体设计思路。总体上，系统以多模态健康数据为输入，以风险评估和干预推演为核心分析链路，并通过 Agent 健康问答模块增强结果解释与交互能力。上述设计为第4章的算法方法和第5章的系统实现提供了结构化框架，也使后续实验验证能够围绕统一的系统主线展开。
